\documentclass[11pt]{article}
\usepackage[hmargin=1in,vmargin=0.86in]{geometry}
\usepackage{amsmath,amssymb,amsthm,booktabs}
\usepackage[hidelinks]{hyperref}
\newtheorem{theorem}{Theorem}
\newtheorem{lemma}{Lemma}
\title{An exhaustive search for small graphic Merino--Welsh counterexamples}
\author{Claude Make No Mistakes}
\date{26 November 2026}
\newcommand{\scope}{33}
\begin{document}
\maketitle

\begin{abstract}
We use the reduction to \(2\)-inflated matroids in the graphic class.
We group parallel core edges
and compute all candidate expansions of each core in a single
exact tensor transform. The search proves the multiplicative
inequality for every loopless bridgeless multigraph with at most
\scope{} edges.
\end{abstract}

\section{Reduction to paired expansions}
Write \(A(G)=T_G(2,0)\), \(C(G)=T_G(0,2)\), and
\(\tau(G)=T_G(1,1)\). For a connected graph these count acyclic
orientations, totally cyclic orientations, and spanning trees.
The multiplicative Merino--Welsh conjecture is
\(A(G)C(G)\geq\tau(G)^2\) for loopless bridgeless \(G\).
The matroid analogue is false \cite{beke}; the graph question
remains open \cite{vali}.

\begin{lemma}
If the graph inequality fails, a minimum-edge counterexample can be
constructed from a connected, loopless, bridgeless multigraph \(H\)
by doubling or subdividing every edge exactly once. In each
parallel bundle of \(H\), at most one edge is doubled.
\end{lemma}
\begin{proof}
Apply Knauer, Mart{\'\i}nez-Sandoval, and Merino's Lemma 1
\cite{kmm} to graphic matroids with at most the minimum
counterexample size. This class is minor-closed. The reduction
gives a counterexample with its elements partitioned into
two-element parallel and series classes. It is connected, since
the three Tutte evaluations multiply over the components and a
counterexample component would be smaller. The paired core
construction \cite[\S3]{kmm}, also given by Chen, Jin, and Ma
\cite{cjm}, yields a connected loopless coloopless graphic
matroid by keeping one representative from each pair. The original
matroid is obtained by one parallel or series extension of each
representative. The corresponding
graph operations are doubling and subdivision. After removing
isolated vertices a graph representing the core has one connected
component and no cut vertex. A cut vertex would partition its
edges so that no circuit meets two parts. Consequently its simple
support is biconnected when it has at least three vertices.
The two-vertex case consists of parallel edges.
Doubling two edges with identical endpoints would create a
parallel class of size at least four in the counterexample,
contrary to the two-element class property.
\end{proof}

\newpage
\section{Exact evaluations}
Let \(m\) and \(n\) be the edge count and order of \(H\). Let \(S\)
be the edges chosen for subdivision and \(P=E(H)\setminus S\)
the edges chosen for doubling. For a subset of expanded edges,
declare a core edge conductive when its doubled pair is nonempty
or its subdivided pair is full. Let \(D\) denote the conductive set
and \(k(D)\) its number of components on the original vertex set.
Components on the expanded vertex set are exactly the components
of \(D\) plus isolated subdivision vertices from empty pairs.
The following expressions for the expanded graph \(G\) are the
graphic specialization of \cite[Proposition 3]{kmm}:
\begin{align*}
 A(G) &= \sum_{D\subseteq E(H)}(-1)^{n+k(D)+|D|}\,3^{|S\setminus D|},\\
 C(G) &= \sum_{D\subseteq E(H)}(-1)^{k(D)-1}\,3^{|P\cap D|},\\
 \tau(G)&=\sum_{\substack{D\subseteq E(H)\\D\text{ a spanning tree}}}
                         2^{|S\setminus D|+|P\cap D|}.
\end{align*}
At \((2,0)\), the Tutte sign is \((-1)^{|F|-|V(G)|+k_G(F)}\)
for an expanded edge set \(F\). The two singleton selections and
the double selection of a conductive parallel pair contribute
\(-2+1=-1\). A nonconductive series pair contributes \(-3\)
from its empty and singleton selections, with the empty state
introducing an isolated vertex. Factoring out one minus sign per
series pair yields the displayed formula for \(A\).
At \((0,2)\), the sign is \((-1)^{k_G(F)-1}\). A conductive
parallel pair has three selections, and a nonconductive series
pair contributes \(2-1=1\). A spanning tree uses exactly one
edge of each chosen parallel pair and one edge of each
nonconductive series pair to attach its interior vertex.
These weights are specializations of standard multivariate Tutte
series and parallel identities \cite{sokal}.

\section{Parallel bundles and exhaustive scan}
Let the simple support have \(\ell\) edges. In a bundle of size \(b\)
the two eligible decorations are one doubled edge with \(b-1\)
subdivisions, and \(b\) subdivisions. An interchange of parallel
core edges preserves the graph's three counts. For each bundle,
rows below represent these two decorations, and columns represent
zero or at least one conductive core edge in the bundle:
\[
 A_b=\begin{pmatrix}3^{b-1}&-3^{b-1}\\3^b&2^b-3^b\end{pmatrix},\quad
 C_b=\begin{pmatrix}1&2^{b+1}-1\\1&2^b-1\end{pmatrix},\quad
 L_b=\begin{pmatrix}2^{b-1}&(b+3)2^b/4\\2^b&b2^{b-1}\end{pmatrix}.
\]
For \(A_b\), the second column is the sum over nonempty selections
of conductive edges, absorbing \((-1)^{|D|}\). For \(C_b\),
each doubled edge contributes weight three when conductive.
For \(L_b\), a chosen support edge must have exactly one conductive
core edge, since two would create a cycle. The tensor factors
operate on support-subset vectors with entries
\((-1)^{n+k(D)}\), \((-1)^{k(D)-1}\), and the indicator of a support
spanning tree, respectively. We apply each factor to every pair
of entries that differ in its bundle bit. This gives all
\(2^\ell\) eligible decorations in \(O(\ell 2^\ell)\) time.
The row sums in absolute value for each factor are at most
\(4^b\), so all intermediate integer values fit a signed 64-bit
word for the sizes considered. Products and comparisons are
exact.

We used nauty 2.9.3 \cite{nauty}. For each order, \texttt{geng -C}
enumerates every simple biconnected support in the edge range,
and \texttt{multig} enumerates positive multiplicities up to
support automorphisms. We handle the two-vertex case directly.
We reduce supports by repeatedly deleting vertices of degree
at most two, with a degree-two vertex's neighbors made adjacent.
This preserves any \(K_4\) minor, since a removed vertex cannot
be a singleton branch set of such a minor. If the support reduces
to the empty graph, all its expansions are series--parallel
and satisfy the inequality \cite{noble}. The exact transform
covers the remaining cores.

\begin{theorem}
Every loopless bridgeless multigraph with at most \scope{} edges
satisfies the multiplicative Merino--Welsh inequality.
\end{theorem}
\begin{proof}
A smallest counterexample with at most \scope{} edges would have
a core with at most \(16\) edges by the lemma.
The search covers each of those cores and found no violation.
\end{proof}

The ledger contains per-order and per-size counts, with
\(7{,}082{,}303\) cores in total. The series--parallel theorem
handles \(2{,}209{,}101\) cores and the transform checked
\(69{,}476{,}958{,}848\) decorations of the remaining cores.
The smallest ratio in the computed part is \(121104/42025\).
It occurs for the core \(K_4\) with three edges of a spanning
path subdivided and the other three doubled.
The independent checker uses deletion--contraction on expanded
graphs for control instances and verifies archived spanning-tree
counts by exact Laplacian determinants.

\subsection*{Further bundle pruning}
The replacement tables of Noble and Royle \cite{noble} also
exclude a bundle of at least three edges with one doubled edge
and a fully subdivided bundle of at least four edges from a
minimum counterexample. For the first case the quotient can be
compared quantitatively. In a connected remainder \(Q\), let
\(a_i,c_i,t_i\) be the three evaluations on \(Q\) for \(i=0\)
and on the terminal contraction for \(i=1\). Adding a mixed bundle
of \(k\) core edges gives
\[
\frac{A_k}{2^{k-1}}=(3/2)^{k-1}(a_0+a_1),\qquad
\frac{C_k}{2^{k-1}}=4c_1+2^{1-k}(c_0-c_1),\qquad
\frac{\tau_k}{2^{k-1}}=t_0+\frac{k+3}{2}t_1.
\]
Contraction sends directed cycles to directed closed walks, so
\(c_1\geq c_0\). For \(k\geq2\), the normalized tree count grows
by a factor at most \(6/5\) when \(k\) increases by one. Therefore
the quotient \(A_kC_k/\tau_k^2\) grows by a factor at least
\((3/2)(5/6)^2=25/24\) per added path. The finite search above
also evaluates the larger bundles.

\begin{thebibliography}{9}
\setlength{\itemsep}{0.18em}
\bibitem{beke} C. Beke, G. K. Csáji, P. Csikvári, and S. Pituk.
The Merino--Welsh conjecture is false for matroids.
\emph{Adv. Math.} 446 (2024), 109674.
\href{https://arxiv.org/abs/2311.01932}{arXiv:2311.01932}.
\bibitem{kmm} K. Knauer, L. Mart{\'\i}nez-Sandoval, and C. Merino.
Merino--Welsh inequalities for matroids with controlled lattices of
cyclic flats.
\href{https://arxiv.org/html/2609.38047v1}{arXiv:2609.38047v1}, 2026.
\bibitem{cjm} J. Chen, X. Jin, and T. Ma.
On minimum-size matroid counterexamples to the multiplicative
Merino--Welsh inequality.
\href{https://arxiv.org/html/2610.27706v1\#S5}{arXiv:2610.27706v1}, 2026.
\bibitem{noble} S. D. Noble and G. F. Royle.
The Merino--Welsh conjecture holds for series--parallel graphs.
\emph{Eur. J. Combin.} 38 (2014), 24--35.
\href{https://arxiv.org/abs/1303.6416}{arXiv:1303.6416}.
\bibitem{sokal} A. D. Sokal. The multivariate Tutte polynomial
(alias Potts model) for graphs and matroids.
\href{https://arxiv.org/abs/math/0503607}{arXiv:math/0503607}, 2005.
\bibitem{nauty} B. D. McKay and A. Piperno.
Practical graph isomorphism, II.
\emph{J. Symbolic Comput.} 60 (2014), 94--112.
\href{https://doi.org/10.1016/j.jsc.2013.09.003}{DOI link}.
\bibitem{vali} B. Váli. The Merino--Welsh conjecture for sparse
and dense graphs. \href{https://arxiv.org/abs/2611.15008v1}{arXiv:2611.15008v1}, 2026.
\end{thebibliography}
\end{document}
